% Part:sets-functions-relations
% Chapter: size-of-sets
% Section: introduction

\documentclass[../../../include/open-logic-section]{subfiles}

\begin{document}

\olfileid[fa-IR]{sfr}{siz}{int}

\olsection{مقدمه}

هنگامی که گئورگ کانتور در دههٔ ۱۸۷۰ نظریهٔ مجموعه‌ها را پروراند، یکی
از هدف‌هایش پذیرفتنی‌ساختن اندیشهٔ یک گردایهٔ نامتناهی بود---یک نامتناهیِ
بالفعل، به تعبیر اندیشمندان قرون وسطی. بخش کلیدی این کار، شیوهٔ پرداخت
او به \emph{اندازهٔ} مجموعه‌های گوناگون بود. اگر $a$، $b$ و $c$
همگی متمایز باشند، آنگاه مجموعهٔ $\{a, b, c\}$ به‌طور شهودی
\emph{بزرگ‌تر} از $\{a, b\}$ است. اما مجموعه‌های نامتناهی چطور؟ آیا
همهٔ آن‌ها هم‌اندازه‌اند؟ معلوم می‌شود که چنین نیست.

نخستین اندیشهٔ مهم در اینجا مفهوم برشماری است. می‌توانیم هر مجموعهٔ
متناهی را با فهرست‌کردن همهٔ !!{element}s آن برشماریم. دربارهٔ برخی
مجموعه‌های نامتناهی نیز، اگر اجازه دهیم خود فهرست نامتناهی باشد،
می‌توانیم همهٔ !!{element}s آن‌ها را فهرست کنیم. چنین مجموعه‌هایی
!!{enumerable} نامیده می‌شوند. نتیجهٔ شگفت‌آور کانتور، که تا پایان این
فصل آن را به‌طور کامل خواهیم فهمید، این بود که برخی مجموعه‌های نامتناهی
!!{enumerable} نیستند.

\end{document}
